\chapter{Balancing}

\section{Problema del class imbalance}

Il \textbf{class imbalance} si verifica quando i possibili valori della variabile di interesse non sono distribuiti in modo equilibrato dal punto di vista della numerosità. Per esempio, se in un dataset sono presenti il 5\% di classi \emph{buggy} e il 95\% di classi \emph{non buggy}, il dataset è fortemente sbilanciato.

Nella maggior parte delle applicazioni reali, lo sbilanciamento è una caratteristica naturale del problema. L'evento che si vuole identificare, cioè la classe positiva, è spesso raro.

\begin{notaslide}
Il class imbalance è frequente in applicazioni come \textbf{fraud detection}, \textbf{intrusion detection}, \textbf{risk management}, classificazione di testi e diagnosi o monitoraggio medico.
\end{notaslide}

I classificatori standard tendono a essere dominati dalla classe maggioritaria. Poiché nel training set osservano molte più istanze della maggioranza, imparano a predire soprattutto quella classe. Il risultato è un'alta accuratezza predittiva sulla classe maggioritaria, ma una bassa capacità predittiva sulla classe minoritaria.

Un input bilanciato può aiutare il classificatore perché rende più visibili i pattern della classe minoritaria, mettendo il modello in condizione di riconoscere meglio i positivi.

\section{Effetto dello sbilanciamento sui classificatori}

In presenza di class imbalance, un classificatore può ottenere una \textbf{accuracy} apparentemente alta predicendo quasi sempre la classe maggioritaria. Tuttavia, questo comportamento può essere inutile rispetto all'obiettivo reale del problema.

\begin{notaprof}
Se lo scopo è trovare frodi o diagnosticare una malattia, un modello che predice sempre ``sano'' o ``non frode'' può ottenere un'accuracy elevata quando i positivi sono pochi, ma fallisce proprio nel compito più importante: identificare i casi rari.
\end{notaprof}

Nel contesto della \textbf{defect prediction}, la classe minoritaria è spesso quella delle classi \emph{buggy}. Il problema non è riconoscere le molte classi sane, ma individuare le poche classi realmente difettose.

Se il training set contiene pochi difetti, il modello tenderà a predire pochi difetti anche nel testing set, riducendo soprattutto la capacità di trovare veri positivi.

\section{Undersampling}

L'\textbf{undersampling} è una tecnica di bilanciamento che riduce il numero di istanze della classe maggioritaria, preservando tutte le istanze della classe minoritaria.

L'idea è estrarre un sottoinsieme più piccolo della maggioranza, in modo che la distribuzione tra le classi diventi più equilibrata.

\begin{notaslide}
Nel diagramma delle slide, un'icona a forma di forbice taglia le righe in eccesso dal blocco della \emph{majority class}, riducendolo fino a renderlo confrontabile con il blocco della \emph{minority class}.
\end{notaslide}

Il vantaggio principale dell'undersampling è la semplicità: il dataset diventa più bilanciato e spesso anche più piccolo, riducendo il costo computazionale dell'addestramento.

Lo svantaggio principale è la \textbf{perdita di informazione}. Eliminando istanze della classe maggioritaria, si possono rimuovere esempi utili per descrivere correttamente il problema.

\begin{notaprof}
L'undersampling può peggiorare le prestazioni del classificatore perché fornisce meno dati al modello. Bilanciare il dataset eliminando esempi significa anche rinunciare a una parte della conoscenza disponibile.
\end{notaprof}

\section{Oversampling}

L'\textbf{oversampling} è una tecnica di bilanciamento che aumenta il numero di istanze della classe minoritaria.

Nella sua forma più semplice, l'oversampling duplica le istanze minoritarie già presenti, fino a rendere la classe minoritaria numericamente confrontabile con la classe maggioritaria.

\begin{notaslide}
Nel diagramma delle slide, una lente di ingrandimento ispeziona il piccolo blocco della \emph{minority class} e lo espande, creando copie fino a raggiungere una dimensione simile a quella della \emph{majority class}.
\end{notaslide}

Il vantaggio principale dell'oversampling è che non elimina dati storici: tutte le istanze originali vengono mantenute.

Lo svantaggio è il rischio di \textbf{overfitting}. Il classificatore può attribuire troppa importanza a istanze duplicate, scambiando ripetizioni artificiali per informazione reale.

\begin{notaprof}
Duplicare dati non significa aggiungere nuova conoscenza. Se un evento raro viene copiato molte volte, il modello può credere che quel pattern sia più rappresentativo di quanto non sia realmente.
\end{notaprof}

\section{SMOTE}

\textbf{SMOTE} significa \textbf{Synthetic Minority Oversampling Technique}. A differenza dell'oversampling semplice, SMOTE non si limita a duplicare istanze esistenti, ma genera nuove istanze sintetiche della classe minoritaria.

La generazione avviene usando le similarità nello \textbf{spazio delle feature} tra istanze minoritarie già presenti. In pratica, SMOTE crea nuovi punti intermedi tra esempi minoritari simili.

\begin{notaslide}
Nel diagramma delle slide, i \emph{majority samples} sono rappresentati come quadrati tratteggiati, i \emph{minority samples} come pallini neri e i \emph{synthetic samples} come pallini rossi. SMOTE genera esempi sintetici posizionandoli lungo segmenti che collegano istanze minoritarie simili.
\end{notaslide}

Il vantaggio di SMOTE è che introduce maggiore variabilità rispetto alla semplice duplicazione. Tuttavia, esiste il rischio che le istanze sintetiche non siano realistiche.

\begin{notaprof}
Il professore sconsiglia l'uso disinvolto di SMOTE. Creare istanze sintetiche significa affidarsi a un algoritmo esterno per inventare nuovi esempi, inserendo dati artificiali nel training set di un altro predittore. Questo può ridurre il realismo dell'esperimento.
\end{notaprof}

\section{Confronto tra undersampling, oversampling e SMOTE}

Le tre tecniche modificano il dataset in modi diversi:

\begin{itemize}
    \item \textbf{undersampling}: riduce la classe maggioritaria;
    \item \textbf{oversampling}: duplica la classe minoritaria;
    \item \textbf{SMOTE}: genera nuove istanze sintetiche della classe minoritaria.
\end{itemize}

L'undersampling può essere utile quando il dataset è molto grande e si vuole ridurre il costo computazionale. L'oversampling può essere utile quando non si vuole perdere informazione storica. SMOTE prova a superare la duplicazione semplice, ma introduce dati artificiali.

I principali rischi sono:

\begin{itemize}
    \item \textbf{perdita di informazione}, tipica dell'undersampling;
    \item \textbf{overfitting}, tipico dell'oversampling semplice;
    \item \textbf{istanze sintetiche non realistiche}, tipiche di SMOTE.
\end{itemize}

\section{Lab}

Il laboratorio richiede di confrontare l'accuratezza dei modelli di machine learning con e senza applicazione dell'oversampling.

Il workflow richiesto è:

\begin{enumerate}
    \item prendere un dataset;
    \item usare tre classificatori:
    \begin{itemize}
        \item \textbf{RandomForest};
        \item \textbf{NaiveBayes};
        \item \textbf{IBk};
    \end{itemize}
    \item applicare una validazione tramite \textbf{66/33 ordered split};
    \item confrontare i risultati con e senza oversampling.
\end{enumerate}

\begin{notaprof}
In un progetto software reale, lo split casuale può essere poco realistico. Se si vuole predire una release futura, il modello deve essere addestrato sulle release precedenti. Mischiare release passate e future può introdurre conoscenza futura nel training, falsando l'esperimento.
\end{notaprof}

\section{Weka API}

\subsection{Undersampling}

\begin{notaslide}
In Weka, l'undersampling può essere eseguito con il filtro:
\[
\texttt{weka.filters.supervised.instance.SpreadSubsample -M 1.0}
\]
\end{notaslide}

\subsection{Oversampling}

\begin{notaslide}
In Weka, l'oversampling può essere eseguito con il filtro:
\[
\texttt{weka.filters.supervised.instance.Resample}
\]
\end{notaslide}

Per l'oversampling, le slide indicano i seguenti parametri:

\begin{itemize}
    \item \texttt{noReplacement=false}, per permettere la duplicazione;
    \item \texttt{biasToUniformClass=1.0};
    \item \texttt{sampleSizePercent=Y}.
\end{itemize}

Il parametro \(Y\) si calcola come:

\[
Y = 100 \cdot \frac{majority - minority}{minority}
\]

\begin{notaslide}
Per il dataset diabetes, le slide riportano l'esempio:
\[
\texttt{weka.filters.supervised.instance.Resample -B 1.0 -Z 130.3}
\]
\end{notaslide}

\subsection{SMOTE}

\begin{notaslide}
In Weka, SMOTE può richiedere l'installazione di un package aggiuntivo tramite il gestore dei pacchetti di Weka.
\end{notaslide}

\section{Collegamento con il progetto}

Il balancing è rilevante nei dataset di defect prediction perché le classi \emph{buggy} sono spesso una minoranza rispetto alle classi \emph{non buggy}.

Senza balancing, un classificatore può imparare a predire quasi sempre \texttt{no}, ottenendo una buona accuracy apparente ma fallendo nel riconoscimento delle classi difettose.

Valutare solo l'accuracy è quindi fuorviante. In un dataset con il 95\% di classi sane, un classificatore che predice sempre \texttt{no} ottiene il 95\% di accuracy, ma non individua alcun difetto.

\begin{notaprof}
Dopo il balancing ci si può aspettare un aumento della \textbf{Recall}, perché il modello riconosce più positivi reali. Tuttavia, può diminuire la \textbf{Precision}, perché aumentando le predizioni positive possono aumentare anche i falsi allarmi.
\end{notaprof}

Per valutare correttamente i modelli, oltre all'accuracy, è opportuno osservare metriche come:

\begin{itemize}
    \item \textbf{Precision};
    \item \textbf{Recall};
    \item \textbf{F1};
    \item \textbf{AUC};
    \item \textbf{Kappa di Cohen}.
\end{itemize}

\section{Da ricordare}

\begin{notaesame}
Il class imbalance è centrale nella defect prediction perché le classi buggy sono spesso rare. Un modello non bilanciato può sembrare accurato, ma ignorare proprio i difetti.
\end{notaesame}

\begin{notaesame}
L'undersampling riduce la classe maggioritaria ma può perdere informazione. L'oversampling mantiene i dati ma può causare overfitting. SMOTE crea dati sintetici, ma questi possono non essere realistici.
\end{notaesame}

\begin{notaesame}
Tutte le attività di bilanciamento devono essere applicate esclusivamente al \textbf{training set}. Non si deve mai bilanciare il testing set. Bilanciare l'intero dataset prima dello split significa falsare la valutazione, perché il classificatore può essere testato su istanze duplicate o artificiali collegate a quelle viste in addestramento.
\end{notaesame}

\begin{notaesame}
In presenza di class imbalance, l'accuracy da sola non basta. È necessario osservare metriche come Recall, Precision, F1, AUC e Kappa.
\end{notaesame}